\chapter{연습 문제 해답}
\label{app:answers}

이 부록은 본문 연습 문제의 해답을 제공한다. 스스로 풀어본 후 확인하기 바란다.

\section{1장 해답}

\textbf{1-1.} ChatGPT의 답변을 평가해보라.

모범 답변: ``언어 모델은 다음 토큰 예측으로 학습되지만, 충분히 큰 모델과 다양한 데이터를 사용하면 질문-답변 패턴을 암묵적으로 학습한다. 사용자가 질문 형태로 입력하면, 모델은 학습 데이터에서 그 질문 다음에 오던 답변 패턴을 생성한다.''

\textbf{1-2.} ``태양은 동쪽에서 ...'' 다음 단어 확률.

사람이라면 ``솟는''을 99\% 이상으로 매긴다. 이유: (1) 일상적 관찰, (2) 언어적 콜로케이션 (``동쪽에서 솟는''은 고정 표현), (3) 문맥의 강한 제약.

\textbf{1-3.} N-gram의 희소성 문제.

``북극곰이 춤을'' 다음 단어를 N-gram이 예측하려면 학습 데이터에 ``북극곰이 춤을'' 다음 단어가 등장했어야 한다. 이 3-gram이 학습되지 않았으면 확률이 0이 되어, smoothing으로 ``더'' 같은 일반적 단어를 추천하게 된다.

\textbf{1-4.} 트랜스포머가 RNN보다 병렬 처리에 유리한 이유.

RNN은 시점 $t$의 계산이 시점 $t-1$의 은닉 상태에 의존하므로 순차 처리 필수. 트랜스포머는 어텐션으로 모든 위치를 동시에 처리하므로 GPU 병렬화에 최적.

\textbf{1-5.} Chinchilla 기준 데이터 스케일.

파라미터를 10배 키우면 데이터도 약 10배 키워야 (20 토큰/파라미터 비율 유지). 즉 10배 모델은 200배 데이터 필요.

\section{3장 해답}

\textbf{3-1.} 빈도가 낮은 쌍을 병합하면?

자주 등장하지 않는 패턴이 토큰이 되어 어휘가 비효율적. 희귀한 ``xz'' 같은 쌍이 병합되면 실제 텍스트에서 거의 사용되지 않는 토큰이 됨.

\textbf{3-2.} Unigram 패널티 \code{-5.0}으로 변경.

어휘에 없는 부분문자열의 패널티가 작아지므로, Viterbi가 더 적극적으로 분할을 시도. 결과적으로 더 많은 작은 토큰으로 분해되는 경향.

\textbf{3-3.} 한국어 BPE, 어휘 100 vs 1000.

어휘 100: ``안녕하세요''가 5개 이상 바이트 토큰으로 분해. 매우 비효율적.
어휘 1000: ``안녕'', ``하세요'' 같은 서브워드 학습 가능. 시퀀스 길이 대폭 단축.

\textbf{3-4.} Viterbi 길이 상한 4로 축소.

5자 이상의 토큰이 사라짐. ``hello'' 같은 단어가 ``he'', ``ll'', ``o''로 분해되어 시퀀스 길이 증가.

\textbf{3-5.} \code{<bos>}의 필요성.

없어도 모델은 동작하지만, 시퀀스 시작을 명시적으로 표시하면 모델이 ``여기서부터 새 시퀀스''임을 인식. 배치 처리 시 여러 시퀀스를 구분하는 데도 유용.

\section{4장 해답}

\textbf{4-1.} 임베딩 차원이 너무 작거나 큰 경우.

작으면 (8): 단어 간 구분이 어려워 의미적 유사성 학습 불가.
크면 (4096): 파라미터 폭발, 메모리 부족, 과적합 위험. 보통 256$\sim$4096 사용.

\textbf{4-2.} Sinusoidal 상수 100 $\to$ 100.

주기가 짧아져 가까운 위치도 구분하기 어려워짐. 위치 인코딩의 ``해상도''가 낮아짐.

\textbf{4-3.} 스케일링 생략.

위치 임베딩이 토큰 임베딩을 압도하여 모델이 위치에만 의존. 의미 정보 손실.

\textbf{4-4.} 학습 가능한 위치 임베딩, max\_len 초과.

인덱스 범위를 벗어나 에러 발생. RoPE 같은 extrapolation이 가능한 방식이 필요.

\textbf{4-5.} LayerNorm $\epsilon$ 값 변경.

$\epsilon$이 크면 정규화가 약해져 출력 분산이 커짐. 너무 작으면 분모가 0에 가까워 수치 불안정.

\section{5장 해답}

\textbf{5-1.} 스케일링 생략 시 softmax.

$d_k=64$일 때 내적 값이 $\pm 64$까지 갈 수 있어 softmax가 원-핫에 포화. 그래디언트 0, 학습 멈춤.

\textbf{5-2.} 인과 마스크 미사용.

학습 시: 모델이 미래 토큰을 ``보고'' 정답을 맞추므로 학습 손실은 낮아지지만 일반화 안 됨.
추론 시: 미래 토큰이 없으므로 학습-추론 갭 발생. 성능 급락.

\textbf{5-3.} 헤드 수 1 (단일 헤드).

하나의 어텐션 패턴만 학습. 문법, 의미, 담화 관계를 한 헤드가 모두 처리해야 하므로 성능 저하.

\textbf{5-4.} 헤드 수 = $d_{model}$ ($d_{head}=1$).

각 헤드가 1차원만 처리. 어텐션 의미가 사라지고 단순 가중합이 됨. 사실상 선형 변환과 동일.

\textbf{5-5.} 어텐션 행렬 메모리, $L=100,000$.

$100000^2 \times 2$ bytes (FP16) = 20 GB. 단일 헤드 기준. 헤드가 32개면 640 GB. 사실상 불가능.

\section{6장 해답}

\textbf{6-1.} $d_{ff} = d_{model}$.

FFN의 확장-축소 패턴이 사라져 비선형 표현력 감소. 더 깊은 선형 변환과 유사.

\textbf{6-2.} 잔차 없이 12층.

backward 시 그래디언트가 앞쪽 층에 도달하지 못해 학습 불가. 손실이 줄어들지 않음.

\textbf{6-3.} Pre-LN에서 \code{ln\_f} 필요성.

Pre-LN은 각 블록 입력을 정규화하지만 마지막 블록 출력은 정규화되지 않은 상태. LM Head에 넣기 전 정규화 필요.

\textbf{6-4.} GELU 대신 SiLU.

성능 차이 미미. SiLU ($x \cdot \sigma(x)$)도 부드러운 전환. 일부 모델(LLaMA의 FFN은 SwiGLU이지만 활성화로 SiLU 사용)이 채택.

\textbf{6-5.} 잔차 스케일링 없이 100층.

초기 활성값 분산이 층마다 누적되어 폭발. 학습 초기에 NaN 발생 가능성 높음.

\section{7장 해답}

\textbf{7-1.} 가중치 타이 미사용, GPT-2 small.

$V \times d = 50257 \times 768 \approx 38.6M$ 파라미터 추가. 총 124M $\to$ 163M.

\textbf{7-2.} Warmup 0.

학습 초기 그래디언트가 불안정하여 발산 가능성. 특히 Post-LN에서 심각. Pre-LN은 덜 민감하지만 여전히 권장.

\textbf{7-3.} Label smoothing 0.3.

강한 smoothing으로 모델이 정답에 대한 확신을 덜 가짐. 일반화는 향상되지만 학습 손실이 높게 유지. 보통 0.1이 적정.

\textbf{7-4.} 그래디언트 클리핑 없음.

그래디언트 폭발 시 가중치가 크게 변하여 모델이 붕괴. 특히 큰 학습률에서 위험.

\textbf{7-5.} PPL 1.0 가능?

이론적으로 불가능. 자연어는 본질적 불확실성이 있어 다음 토큰이 100\% 예측 불가. 학습 데이터에 완전히 외운 경우에만 특정 문장에서 PPL $\approx$ 1 가능하지만 일반적이지 않음.

\section{8장 해답}

\textbf{8-1.} Greedy 반복 루프 원인.

분포가 ``store'' 다음 ``to buy''를 높게, ``to buy'' 다음 ``some store''를 높게 예측하는 루프가 형성되면 greedy는 거기서 빠져나오지 못함. 확률이 높은 경로만 따라가기 때문.

\textbf{8-2.} Temperature $\to$ 0.

$T=0$이면 softmax가 argmax가 되어 greedy와 동일. $T=0.01$이면 거의 greedy. 분포가 너무 날카로워져 다양성 0.

\textbf{8-3.} Top-K 후 Top-P.

Top-K로 후보를 K개로 줄인 후, 그 중에서 누적 확률 P까지 선택. 더 강력한 필터링. 품질은 향상되지만 다양성 감소.

\textbf{8-4.} 챗봇에 적합한 샘플링.

Top-P (P=0.9, T=0.7$\sim$0.9). 다양성과 일관성의 균형. Greedy는 단조로움, 순수 샘플링은 일관성 부족.

\textbf{8-5.} Speculative decoding이 빠른 이유.

작은 모델이 K개 토큰을 빠르게 생성하면, 큰 모델이 그 K개를 한 번에 검증 (병렬 처리). 맞으면 K개를 한 스텝에 채택하므로 큰 모델의 스텝 수가 1/K로 감소.
